\hypertarget{glossuxe1rio-tuxe9cnico-sistemas-scada-e-supervisuxf3rios-industriais}{%
\chapter{\texorpdfstring{Glossário Técnico --- Sistemas SCADA e Supervisórios Industriais}{Glossário Técnico  Sistemas SCADA e Supervisórios Industriais}}\label{glossuxe1rio-tuxe9cnico-sistemas-scada-e-supervisuxf3rios-industriais}}

\begin{caixanota}

\textbf{Como usar:} As siglas são expandidas na primeira ocorrência de cada capítulo. Este glossário reúne todos os termos para consulta rápida.

\end{caixanota}

\begin{center}\rule{0.5\linewidth}{0.5pt}\end{center}

\hypertarget{a}{%
\section{A}\label{a}}

\textbf{Air Gap}\\
Isolamento físico completo entre uma rede (geralmente OT) e outras redes externas. Historicamente usado em sistemas SCADA críticos para impedir acesso remoto --- mas eliminado pela convergência TI/OT.

\textbf{AMQP} (\emph{Advanced Message Queuing Protocol})\\
Protocolo de mensageria orientado a filas, similar ao MQTT mas com mais recursos de roteamento. Usado em middlewares enterprise como RabbitMQ.

\textbf{Analytics Industrial}\\
Conjunto de técnicas de análise de dados aplicadas ao ambiente industrial: estatística descritiva, detecção de anomalias, modelos preditivos.

\textbf{Alarme}\\
Evento que exige \textbf{resposta do operador} dentro de um tempo definido. Se não há ação possível, não é alarme: é registro.

\textbf{Alarme fixo} (\emph{standing alarm})\\
Alarme que permanece ativo por longos períodos sem ação possível, geralmente por falha de instrumento ou de configuração.

\textbf{Alarme obsoleto} (\emph{stale alarm})\\
Alarme cuja condição já cessou mas que permanece ativo por falha de projeto ou de comunicação.

\textbf{Alarme fugaz} (\emph{fleeting alarm})\\
Alarme que dispara e limpa em poucos segundos, normalmente por transitório de partida; combate-se com atraso de confirmação.

\textbf{Alarm flooding (inundação de alarmes)}\\
Rajada de alarmes em curto intervalo que ultrapassa a capacidade de leitura e de resposta do operador, escondendo a causa raiz.

\textbf{Anomalia}\\
Comportamento estatisticamente diferente do padrão observado. Não é sinônimo de falha: pode ser regime novo, produto diferente ou condição ambiental. Serve como convite à investigação, não como alarme.

\textbf{Atributo (ThingsBoard)}\\
Valor de estado associado a uma entidade, em que o último valor substitui o anterior (ao contrário da telemetria, que preserva o histórico). Escopos: servidor, compartilhado e cliente.

\textbf{Atuador}\\
Dispositivo que converte um sinal elétrico em ação física sobre o processo: válvulas, motores, relés, pistões pneumáticos.

\begin{center}\rule{0.5\linewidth}{0.5pt}\end{center}

\hypertarget{b}{%
\section{B}\label{b}}

\textbf{Baseline (linha de base)}\\
Referência de comportamento normal de um ativo ou processo, obtida do próprio histórico, com média, dispersão e correlações. Sem baseline, não há como afirmar que um valor está anormal.

\textbf{BPCS} (\emph{Basic Process Control System})\\
Sistema de controle principal de uma planta --- geralmente o DCS ou CLP. Distinto do SIS (sistema de segurança).

\textbf{Broker MQTT}\\
Servidor intermediário no padrão MQTT que recebe mensagens dos publishers e as distribui aos subscribers inscritos nos tópicos correspondentes.

\begin{center}\rule{0.5\linewidth}{0.5pt}\end{center}

\hypertarget{c}{%
\section{C}\label{c}}

\textbf{CapEx} (\emph{Capital Expenditure})\\
Investimento de capital em ativos físicos duráveis --- servidores, hardware, licenças perpétuas.

\textbf{Chattering (alarme repetitivo)}\\
Alarme que dispara e limpa repetidamente em torno do limite, consumindo a atenção do operador; combate-se com faixa morta e atraso de confirmação.

\textbf{CIA} (\emph{Confidentiality, Integrity, Availability})\\
Tríade clássica de segurança da informação: Confidencialidade, Integridade e Disponibilidade.

\textbf{CIP} (\emph{Common Industrial Protocol})\\
Protocolo de aplicação usado pelo EtherNet/IP e DeviceNet para comunicação industrial Rockwell/ODVA.

\textbf{Ciclo de vida de um projeto SCADA}\\
Sequência de fases --- definição, seleção, projeto, implementação, testes preliminares, aceitação e manutenção --- em que cada fase termina com um produto verificável.

\textbf{CLP} (\emph{Controlador Lógico Programável})\\
Computador industrial programável para controle de sequências lógicas e processos discretos. Equivalente inglês: PLC (\emph{Programmable Logic Controller}).

\textbf{Cloud SCADA}\\
Arquitetura em que o software SCADA é hospedado em infraestrutura de nuvem (AWS, Azure, GCP) em vez de servidores físicos na planta.

\textbf{CoAP} (\emph{Constrained Application Protocol})\\
Protocolo leve para dispositivos IoT com recursos limitados, similar ao HTTP mas baseado em UDP. Ideal para sensores com baixa energia.

\textbf{Compressão por exceção (swinging door)}\\
Estratégia de historização que só grava um novo ponto quando o valor se afasta do último gravado mais que uma tolerância definida. A tolerância é o erro máximo da série reconstruída, e o descarte é irreversível.

\begin{center}\rule{0.5\linewidth}{0.5pt}\end{center}

\hypertarget{d}{%
\section{D}\label{d}}

\textbf{Contêiner}\\
Forma de empacotar uma aplicação com suas dependências em uma imagem versionada, executável de forma isolada --- usada na borda para atualizar aquisição, regras e modelos sem trocar o sistema do gateway.

\textbf{Dashboard}\\
Painel que reúne indicadores, tendências, estados e alarmes para apoiar decisão. Sem contexto (faixa normal, limite e período) o valor exibido não informa.

\textbf{DCS} (\emph{Distributed Control System} --- Sistema de Controle Distribuído*)\\
Sistema de controle para processos contínuos de grande escala, com controle distribuído em múltiplos controladores de campo e banco de dados centralizado.

\textbf{Deadband (faixa morta)}\\
Diferença necessária entre o valor de disparo e o de limpeza de um alarme, evitando que ele oscile em torno do limite.

\textbf{Deriva (drift)}\\
Perda de validade de um modelo ou de um baseline em produção, causada por mudança do processo, do sensor ou da relação entre as variáveis. Silenciosa: o sistema continua respondendo, com confiança, errado.

\textbf{Digital Twin (Gêmeo Digital)}\\
Representação virtual dinâmica de um ativo físico, processo ou sistema que sincroniza em tempo real com seu equivalente real, permitindo simulação e otimização.

\textbf{DMZ Industrial}\\
\emph{Demilitarized Zone} --- zona de rede intermediária entre a rede corporativa (TI) e a rede industrial (OT), contendo sistemas que precisam trocar dados com ambos os lados (ex.: historiador, servidor web).

\textbf{DNP3} (\emph{Distributed Network Protocol 3})\\
Protocolo de comunicação para sistemas SCADA de infraestrutura crítica (energia, água, óleo e gás). Suporta reporte espontâneo e timestamping de eventos.

\textbf{DRP} (\emph{Disaster Recovery Plan})\\
Plano de recuperação de desastres --- conjunto de procedimentos para restaurar sistemas críticos após falha grave.

\begin{center}\rule{0.5\linewidth}{0.5pt}\end{center}

\hypertarget{e}{%
\section{E}\label{e}}

\textbf{Edge Computing (Computação na Borda)}\\
Processamento de dados realizado próximo à fonte --- no próprio dispositivo ou em um gateway de campo --- em vez de na nuvem central.

\textbf{EEMUA 191}\\
Guia de engenharia da EEMUA para projeto, gestão e contratação de sistemas de alarme; 3ª edição, 2013. É guia, não norma, e é referência usual de valores de desempenho.

\textbf{Escalonamento (de faixa)}\\
Conversão de uma leitura bruta (contagem do conversor A/D) para a unidade de engenharia do instrumento, por regra de três entre os extremos da faixa.

\textbf{Espaço de endereços (address space)}\\
No OPC UA, a estrutura de nós e referências que o servidor expõe e que o cliente descobre com \passthrough{\lstinline!Browse!} --- diferente do mapa fixo de registradores de um protocolo clássico.

\textbf{Ergonomia}\\
Conjunto de disciplinas que estuda a organização do trabalho na interação entre seres humanos e máquinas. Aplicada à IHM, busca reduzir sobrecarga, monotonia, cansaço e excesso de informação.

\textbf{ERP} (\emph{Enterprise Resource Planning} --- Planejamento de Recursos Empresariais*)\\
Sistema de gestão integrada de negócios: finanças, RH, logística, vendas. Integra-se ao MES via ISA-95.

\textbf{Estado do equipamento}\\
Condição qualitativa de um ativo (parado, partindo, operando, parando, em falha), modelada explicitamente em vez de deduzida na tela a partir de tags discretos.

\textbf{EtherNet/IP}\\
Protocolo industrial Ethernet baseado no CIP (\emph{Common Industrial Protocol}), desenvolvido pela ODVA e amplamente usado no ecossistema Rockwell/Allen-Bradley.

\textbf{Evento}\\
Ocorrência com importância para o operador que \textbf{não} exige ação corretiva (partida de equipamento, troca de turno, janela aberta). Distingue-se do alarme, que exige intervenção.

\textbf{Exatidão}\\
Grau de proximidade entre o valor medido e o verdadeiro. É definida pelo transmissor e pela instalação --- não pela resolução do conversor A/D, que apenas determina o tamanho do degrau.

\begin{center}\rule{0.5\linewidth}{0.5pt}\end{center}

\hypertarget{f}{%
\section{F}\label{f}}

\textbf{Failover}\\
Comutação automática de um sistema primário para um backup em caso de falha, sem intervenção humana.

\textbf{FFT} (\emph{Fast Fourier Transform})\\
Algoritmo de transformada de Fourier para converter um sinal no domínio do tempo para o domínio da frequência --- essencial para análise de vibração em manutenção preditiva.

\textbf{Fieldbus}\\
Barramento digital de comunicação para interligar sensores, atuadores e controladores no nível de campo. Exemplos: Profibus, Foundation Fieldbus, DeviceNet, AS-i.

\textbf{Fog Computing}\\
Camada de computação intermediária entre o edge (campo) e a nuvem --- geralmente um servidor local na planta ou subestação que agrega e pré-processa dados.

\textbf{Foundation Fieldbus (FF)}\\
Protocolo de fieldbus digital para instrumentação de processo contínuo (petroquímica, refino), com capacidade de executar blocos de função (controle) no próprio instrumento.

\begin{center}\rule{0.5\linewidth}{0.5pt}\end{center}

\hypertarget{g}{%
\section{G}\label{g}}

\textbf{Gateway de Protocolo}\\
Dispositivo ou software que converte entre diferentes protocolos de comunicação (ex.: Modbus RTU \ensuremath{\rightarrow} OPC UA).

\textbf{Gateway de borda}\\
Equipamento que materializa a borda de campo: conversa com o processo pela face sul (Modbus, OPC UA, DNP3), normaliza e publica pela face norte (MQTT, REST, Sparkplug B), com fila local para sobreviver à queda do enlace.

\textbf{Grafana}\\
Plataforma de visualização e monitoramento que consulta fontes externas (séries temporais, bancos SQL, Prometheus). Não armazena o dado de processo e não é sistema de comando.

\begin{center}\rule{0.5\linewidth}{0.5pt}\end{center}

\hypertarget{h}{%
\section{H}\label{h}}

\textbf{Hierarquia de navegação}\\
Organização das telas de um supervisório do geral para o particular, com barra de navegação fixa, navegação horizontal pelo fluxo do processo e retorno sempre disponível.

\textbf{HART} (\emph{Highway Addressable Remote Transducer})\\
Protocolo que superpõe comunicação digital (FSK --- Frequency Shift Keying) sobre o sinal analógico 4-20 mA, permitindo configuração e diagnóstico remoto de instrumentos.

\textbf{Historiador (historian)}\\
Sistema que armazena séries temporais de processo, com compressão por exceção e consulta por intervalo. Distingue-se do histórico de eventos, que registra ocorrências com estado, prioridade e reconhecimento.

\textbf{HMI} (\emph{Human-Machine Interface} --- Interface Homem-Máquina*)\\
Dispositivo ou software que apresenta o estado do processo ao operador e permite intervenções. Pode ser local (painel touchscreen) ou centralizado (estação SCADA).

\textbf{HA} (\emph{High Availability} --- Alta Disponibilidade*)\\
Arquitetura de sistema projetada para minimizar o tempo de inatividade, tipicamente com redundância ativa/passiva ou ativa/ativa.

\begin{center}\rule{0.5\linewidth}{0.5pt}\end{center}

\hypertarget{i}{%
\section{I}\label{i}}

\textbf{IEC 61131-3}\\
Norma internacional que define cinco linguagens de programação para CLPs: Ladder (LD), Blocos de Função (FBD), Texto Estruturado (ST), Lista de Instruções (IL) e Diagrama de Função Sequencial (SFC).

\textbf{IEC 62443}\\
Família de normas internacionais para cibersegurança em sistemas de automação e controle industrial (IACS).

\textbf{IEC 62682}\\
Norma internacional de gerenciamento de alarmes para indústrias de processo, alinhada à ISA-18.2; edição vigente de 2022, que substitui a de 2014.

\textbf{IHM de alta performance} (\emph{High Performance HMI})\\
Estilo de projeto de IHM difundido pela ISA-101: fundo neutro, cor reservada a alarme e anomalia, hierarquia de telas navegável, mínimo de elementos decorativos e foco em consciência situacional e tempo de resposta.

\textbf{IIoT} (\emph{Industrial Internet of Things} --- Internet Industrial das Coisas*)\\
Aplicação do paradigma IoT ao ambiente industrial, com requisitos rigorosos de confiabilidade, latência e segurança.

\textbf{InfluxDB}\\
Banco de dados de séries temporais muito usado com o Grafana. O dado medido entra como \emph{field} e o ativo como \emph{tag} --- inverter isso destrói o desempenho.

\textbf{ISA-18.2}\\
Norma ISA para gerenciamento de sistemas de alarmes em indústrias de processo; edição vigente de 2016 (a primeira é de 2009).

\textbf{ISA-88}\\
Norma ISA para controle de processos em batelada (\emph{batch control}): define modelos de receita, equipamentos e fases de processo.

\textbf{ISA-95}\\
Norma ISA para integração entre sistemas de controle industrial e sistemas corporativos (MES e ERP).

\textbf{ISA-101}\\
Norma ISA de projeto de interfaces homem-máquina (IHM) para sistemas de automação de processo --- principal referência global de \textbf{IHM/SCADA de alta performance}: foco cognitivo, uso estratégico de cores, hierarquia da informação e decisão rápida. \textbf{Não é um software SCADA}, e sim um conjunto de diretrizes de engenharia para telas eficientes.

\begin{center}\rule{0.5\linewidth}{0.5pt}\end{center}

\hypertarget{jl}{%
\section{\texorpdfstring{J--L}{JL}}\label{jl}}

\textbf{Jump Server (Bastion Host)}\\
Servidor intermediário através do qual todos os acessos remotos à rede OT devem passar --- nunca acesso direto aos CLPs/SCADA.

\textbf{KPI} (\emph{Key Performance Indicator})\\
Indicador-chave de desempenho. Em dashboard industrial, exige fórmula e período explícitos --- caso contrário vira número decorativo.

\textbf{Ladder (LD)}\\
Linguagem de programação de CLP que representa a lógica de controle em forma de diagrama elétrico com contatos e bobinas. É a linguagem mais utilizada na indústria.

\textbf{LWT} (\emph{Last Will and Testament})\\
Mecanismo MQTT em que o broker publica automaticamente uma mensagem predefinida quando um cliente desconecta inesperadamente --- usado para detecção de falha de dispositivos IIoT.

\begin{center}\rule{0.5\linewidth}{0.5pt}\end{center}

\hypertarget{m}{%
\section{M}\label{m}}

\textbf{Máquina de estados}\\
Modelo que descreve os estados possíveis de um equipamento ou de um procedimento e as transições permitidas entre eles, com o evento que provoca cada transição.

\textbf{Mapa de memória (tabela de alocação)}\\
Documento que relaciona cada tagname do supervisório ao endereço real do hardware de aquisição/controle e ao driver de comunicação responsável pela leitura.

\textbf{Mapa de memória (tabela de alocação)}\\
Documento que relaciona cada tagname do supervisório ao endereço real do hardware de aquisição/controle e ao driver de comunicação responsável pela leitura.

\textbf{Manutenção preditiva}\\
Nível de manutenção em que a intervenção é programada a partir da projeção de falha em janela útil. Exige histórico longo com eventos rotulados; a baseada em condição registra o desvio e o limite que o caracteriza.

\textbf{Média móvel}\\
Filtro digital que suaviza o ruído calculando a média das \emph{n} últimas amostras; reduz o ruído na proporção da raiz quadrada de \emph{n}, ao custo de atrasar a resposta do sinal.

\textbf{MES} (\emph{Manufacturing Execution System} --- Sistema de Execução da Manufatura*)\\
Sistema que conecta o chão de fábrica ao ERP: rastreabilidade de lotes, qualidade, sequenciamento de ordens, OEE.

\textbf{MFA} (\emph{Multi-Factor Authentication})\\
Autenticação multifator --- exige dois ou mais fatores de verificação (senha + token + biometria) para acesso a sistemas críticos.

\textbf{Modbus}\\
Protocolo de comunicação serial criado pela Modicon em 1979. Disponível nas variantes RTU (serial), ASCII e TCP/IP (Ethernet). O mais difundido na indústria.

\textbf{MPC} (\emph{Model Predictive Control})\\
Estratégia de controle avançado que utiliza um modelo matemático do processo para prever comportamento futuro e calcular a ação de controle ótima considerando restrições.

\textbf{MTU} (\emph{Master Terminal Unit})\\
Servidor central de um sistema SCADA responsável pelo polling das RTUs, banco de dados em tempo real, processamento de alarmes e interface com os clientes HMI.

\textbf{MQTT} (\emph{Message Queuing Telemetry Transport})\\
Protocolo leve de mensageria baseado em publicação/assinatura (pub/sub), amplamente utilizado em IIoT. Opera sobre TCP/IP na porta 1883 (ou 8883 com TLS).

\begin{center}\rule{0.5\linewidth}{0.5pt}\end{center}

\hypertarget{n}{%
\section{N}\label{n}}

\textbf{NAMUR NA 102}\\
Recomendação da NAMUR (2003) para gerenciamento de alarmes na indústria química e farmacêutica; referência histórica que influenciou ISA-18.2 e IEC 62682.

\textbf{Node-RED}\\
Plataforma open-source de programação visual \emph{flow-based} para integração de dispositivos, APIs e serviços. Mantida pela OpenJS Foundation.

\textbf{NTP} (\emph{Network Time Protocol})\\
Protocolo de sincronização de relógios em rede, com precisão típica de milissegundos. Pré-requisito para comparar carimbos de tempo entre controlador, supervisório e historian.

\begin{center}\rule{0.5\linewidth}{0.5pt}\end{center}

\hypertarget{o}{%
\section{O}\label{o}}

\textbf{OEE} (\emph{Overall Equipment Effectiveness})\\
Indicador composto de eficiência global de equipamentos = Disponibilidade \ensuremath{\times} Desempenho \ensuremath{\times} Qualidade. Padrão SEMI E10.

\textbf{OpEx} (\emph{Operational Expenditure})\\
Despesas operacionais recorrentes --- energia, manutenção, assinaturas de software, salários.

\textbf{OPC} (\emph{OLE for Process Control})\\
Padrão de interoperabilidade entre softwares SCADA e CLPs no ambiente Windows, baseado em COM/DCOM. Predecessor do OPC UA.

\textbf{OPC UA} (\emph{OPC Unified Architecture})\\
Versão moderna e plataforma-independente do OPC: segurança nativa (TLS), modelo de informação orientado a objetos, suporte a Linux e embarcados. Porta padrão: 4840.

\textbf{OT} (\emph{Operational Technology})\\
Hardware e software que detecta ou causa mudanças por meio do monitoramento e/ou controle direto de equipamentos físicos, processos e eventos industriais. Contrasta com TI (\emph{Information Technology}).

\begin{center}\rule{0.5\linewidth}{0.5pt}\end{center}

\hypertarget{pq}{%
\section{\texorpdfstring{P--Q}{PQ}}\label{pq}}

\textbf{PAM} (\emph{Privileged Access Management})\\
Solução de gestão de acessos privilegiados --- registra, audita e controla sessões de usuários com permissões elevadas em sistemas críticos.

\textbf{PID} (\emph{Proportional-Integral-Derivative})\\
Algoritmo de controle realimentado amplamente utilizado na indústria para manter uma variável de processo no setpoint desejado.

\textbf{Polling}\\
Método de comunicação em que o mestre (MTU/SCADA) interroga periodicamente cada escravo (RTU/CLP) para coletar dados --- contrasta com o reporte espontâneo do DNP3.

\textbf{Prioridade de alarme}\\
Urgência da resposta do operador, derivada da consequência e do tempo disponível. Não confundir com severidade, que descreve a consequência da condição anormal.

\textbf{PROFINET}\\
Protocolo Ethernet industrial da Siemens/PI (PROFIBUS \& PROFINET International), sucessor do Profibus. Suporta RT (Real-Time) e IRT (Isochronous Real-Time) para motion control.

\textbf{Profibus}\\
Protocolo de fieldbus serial criado na Alemanha em 1987. Variantes: DP (automação discreta, até 12 Mbps) e PA (instrumentação de processo em zonas de risco).

\textbf{Provisionamento (as code)}\\
Configuração de dashboards, fontes de dados e alertas por arquivo versionado, em vez de ajustes manuais pela interface --- o que dá auditabilidade e repetibilidade.

\textbf{PTP} (\emph{Precision Time Protocol}, IEEE 1588)\\
Protocolo de sincronização com precisão de sub-microssegundo em rede local, usado quando o carimbo de tempo na origem precisa reconstruir a sequência de eventos.

\textbf{Qualidade do dado (StatusCode)}\\
Código associado a cada valor lido que informa se a leitura é válida (\passthrough{\lstinline!Good!}), incerta (\passthrough{\lstinline!Uncertain!}) ou inválida (\passthrough{\lstinline!Bad!}). Deve ser exibido na tela em vez de substituído por zero.

\textbf{QoS} (\emph{Quality of Service})\\
Nível de garantia de entrega no MQTT: 0 (at-most-once), 1 (at-least-once), 2 (exactly-once).

\begin{center}\rule{0.5\linewidth}{0.5pt}\end{center}

\hypertarget{r}{%
\section{R}\label{r}}

\textbf{Racionalização de alarmes}\\
Processo de documentar, alarme por alarme, causa, consequência, ação esperada, tempo de resposta e prioridade. Alarme sem ação esperada não foi racionalizado.

\textbf{Registrador de primeira falha (RPE)}\\
Lógica no controlador que grava o código do primeiro sensor a atuar em um endereço de memória e bloqueia as gravações seguintes, preservando a causa raiz em meio a uma avalanche de eventos.

\textbf{Resolução de entrada analógica}\\
Menor degrau que o conversor A/D consegue representar, dado pelo número de bits da faixa. Não deve ser confundida com exatidão nem com ruído.

\textbf{Retenção (política de)}\\
Conjunto de regras que define por quanto tempo cada classe de dado é mantida e com que resolução, prevendo agregação, cópia de segurança e descarte planejado.

\textbf{Retenção (política de)}\\
Conjunto de regras que define por quanto tempo cada classe de dado é mantida e com que resolução, prevendo agregação, cópia de segurança e descarte planejado.

\textbf{Rotulagem}\\
Registro da causa e do componente em cada ocorrência de manutenção. É o que transforma o histórico de operação em dado utilizável para aprendizado supervisionado.

\textbf{RPC} (\emph{Remote Procedure Call})\\
Chamada de procedimento remoto. Em plataformas IIoT, permite que a plataforma envie comandos ao dispositivo (lado do servidor) ou que o dispositivo solicite algo à plataforma (lado do dispositivo).

\textbf{RTDB} (\emph{Real-Time Database})\\
Banco de dados em tempo real --- mantido na memória do servidor SCADA com o estado atual de todas as variáveis do sistema.

\textbf{RTU} (\emph{Remote Terminal Unit})\\
Dispositivo de campo que coleta dados de sensores, executa lógicas simples e transmite ao servidor SCADA (MTU) via protocolo de comunicação.

\textbf{Rule chain (cadeia de regras)}\\
No ThingsBoard, sequência visual de nós por onde cada mensagem passa: persistir telemetria, salvar atributos, criar e limpar alarmes, publicar em sistemas externos.

\begin{center}\rule{0.5\linewidth}{0.5pt}\end{center}

\hypertarget{s}{%
\section{S}\label{s}}

\textbf{SCADA} (\emph{Supervisory Control and Data Acquisition})\\
Sistema de controle supervisório e aquisição de dados --- monitora e controla remotamente processos industriais e infraestruturas críticas a partir de uma central de operações.

\textbf{Sensor}\\
Dispositivo que mede uma grandeza física do processo (temperatura, pressão, nível, vazão, posição) e a converte em sinal elétrico padronizado (4-20 mA, 0-10 V, pulso, digital).

\textbf{Sensor virtual (soft sensor)}\\
Modelo que estima em tempo real uma variável medida apenas em laboratório ou por analisador caro, a partir de variáveis medidas continuamente. Precisa declarar a sua região de validade.

\textbf{Série temporal}\\
Conjunto de medições indexadas no tempo. Exige carimbo de tempo confiável: sem o \emph{timestamp} do dispositivo, dado atrasado é gravado como atual.

\textbf{Severidade}\\
Consequência da condição anormal se ninguém agir. Dimensão distinta da prioridade, que trata da urgência da resposta.

\textbf{SFC} (\emph{Sequential Function Chart})\\
Linguagem de programação de CLP (IEC 61131-3) baseada em estados e transições --- ideal para controle de sequências e processos batch.

\textbf{Shelving (silenciamento temporário)}\\
Ocultação manual de um alarme por prazo definido e com justificativa registrada. Diferente da supressão, que é definida em projeto.

\textbf{SIS} (\emph{Safety Instrumented System})\\
Sistema de segurança funcional que monitora condições críticas do processo e executa ações de proteção automática (shutdown) para evitar acidentes. Normatizado pela IEC 61511.

\textbf{SL} (\emph{Security Level})\\
Nível de segurança definido pela IEC 62443: SL1 (proteção contra causas não intencionais) a SL4 (proteção contra atores estatais).

\textbf{SOE} (\emph{Sequence of Events})\\
Registro de sequência de eventos com carimbo de tempo na origem, com resolução de milissegundo ou melhor. Exige relógios sincronizados (NTP ou PTP) em todos os nós.

\textbf{Sparkplug B}\\
Especificação sobre MQTT com esquema de dados padronizado para IIoT industrial --- define payloads, timestamps, estados de sessão e gerenciamento de ativos.

\textbf{StatusCode (OPC UA)}\\
Código que acompanha cada valor lido, indicando se a leitura é válida ou se a qualidade está comprometida. Ignorá-lo faz um sensor com defeito aparecer como ``valor zero''.

\textbf{Store and forward}\\
Armazenamento local do dado durante a queda do enlace, com reenvio posterior preservando o carimbo de origem. É o que evita buracos na série temporal justamente nos períodos de problema.

\textbf{Supressão (suppression)}\\
Ocultação temporária de um alarme por projeto, enquanto a causa que o tornaria irrelevante está ativa (por exemplo, durante o enchimento de um tanque).

\begin{center}\rule{0.5\linewidth}{0.5pt}\end{center}

\hypertarget{t}{%
\section{T}\label{t}}

\textbf{Tag}\\
Identificador único de uma variável de processo em um sistema SCADA ou historiador (ex.: \passthrough{\lstinline!PI-1001!} para pressão no ponto 1001). É o nome pelo qual CLP, supervisório, histórico, relatórios e nuvem se referem ao mesmo dado.

\textbf{Tag de I/O (de comunicação)}\\
Tagname cujo valor é atualizado pelo servidor de comunicação (\emph{driver}), que lê e escreve a memória do hardware de controle. Exige endereçamento informado no cadastro.

\textbf{Tag interno (de memória ou RAM)}\\
Tagname cujo valor é atualizado no próprio software de supervisão --- cálculos, contadores e estado da aplicação ---, sem vínculo de atualização com outro software.

\textbf{Taxa de amostragem}\\
Intervalo entre duas amostras consecutivas de uma variável. Deve seguir a dinâmica do processo: subamostrar produz uma curva que não corresponde à variação real.

\textbf{Taxa de varredura (scan rate)}\\
Frequência com que o mestre interroga os escravos em uma comunicação por \emph{polling}; define o atraso máximo entre a mudança no campo e sua chegada ao supervisório.

\textbf{Telemetria}\\
Série temporal publicada por um dispositivo (temperatura, pressão, vazão): cada valor é um ponto no tempo e o histórico é preservado.

\textbf{Tendência real e tendência histórica}\\
A real exibe os últimos valores a partir de um vetor em memória, com janela limitada; a histórica recupera valores gravados em disco e responde por qualquer período retido.

\textbf{ThingsBoard}\\
Plataforma de código aberto para coleta, processamento, visualização e gestão de dados de dispositivos IoT/IIoT, com motor de regras, alarmes e RPC.

\textbf{TI} (\emph{Tecnologia da Informação})\\
Sistemas computacionais para gestão de informações empresariais: email, ERP, bancos de dados corporativos. Contrasta com OT.

\textbf{TLS} (\emph{Transport Layer Security})\\
Protocolo de criptografia para comunicações seguras em rede --- sucessor do SSL. Usado para proteger conexões OPC UA, MQTT e HTTPS.

\textbf{TSDB} (\emph{Time Series Database})\\
Banco de dados otimizado para séries temporais: índice temporal nativo, compressão por série, agregação contínua e política de retenção própria.

\begin{center}\rule{0.5\linewidth}{0.5pt}\end{center}

\hypertarget{vz}{%
\section{\texorpdfstring{V--Z}{VZ}}\label{vz}}

\textbf{Unidade de engenharia (EU)}\\
Grandeza física em que a variável é apresentada ao operador --- graus Celsius, metros cúbicos por hora, percentual ---, resultante do escalonamento da faixa bruta do conversor.

\textbf{Validação de faixa}\\
Tratamento que rejeita leituras fisicamente impossíveis (fora da faixa do sensor) e as marca com qualidade ruim, em vez de apresentá-las como valor válido.

\textbf{VPN} (\emph{Virtual Private Network})\\
Rede privada virtual que cria um túnel criptografado sobre a internet pública para comunicação segura entre sites remotos ou usuários e a rede corporativa/industrial.

\textbf{WirelessHART}\\
Extensão sem fio do protocolo HART, baseada em IEEE 802.15.4 e malha de rádio (\emph{mesh}). Permite adicionar comunicação sem fio a instrumentos de processo industriais.

\textbf{Zero Trust}\\
Modelo de segurança que pressupõe que nenhuma entidade (usuário, dispositivo ou rede) é confiável por padrão --- toda conexão deve ser autenticada e autorizada explicitamente.
